\documentclass[11pt]{article}

\usepackage{amsmath,amssymb,amsthm}
\usepackage[margin=1in]{geometry}
\usepackage{hyperref}
\usepackage{listings}
\usepackage{xcolor}
\usepackage{microtype}

\hypersetup{
  colorlinks=true,
  linkcolor=blue,
  citecolor=blue,
  urlcolor=blue
}

\lstdefinestyle{julia}{
  basicstyle=\ttfamily\small,
  breaklines=true,
  frame=single,
  columns=fullflexible,
  showstringspaces=false,
  keepspaces=true,
  numbers=none
}
\lstset{style=julia}

\newtheorem{theorem}{Theorem}
\newtheorem{definition}{Definition}
\newtheorem{example}{Example}
\newtheorem{remark}{Remark}

\newcommand{\Ann}{\operatorname{Ann}}
\newcommand{\Rad}{\operatorname{Rad}}
\newcommand{\Ass}{\operatorname{Ass}}

\title{A Computational Verification of Annihilator-Based\\
Radical Dependency via OSCAR}

\author{Hamdullah \"Ozkaya\\
\small Department of Mathematics, Bursa Technical University, Bursa, Turkey\\
\small \texttt{hamdullah.ozkaya@btu.edu.tr}}

\date{\today}

\begin{document}
\maketitle

\begin{abstract}
This note provides an independent computational verification, using the
OSCAR computer algebra system, of six theorems from a recent paper on
annihilator-based dependency relations for submodules of a Noetherian
module. We reproduce the paper's elementary examples over
$\mathbb{Z}/n\mathbb{Z}$, test the characterization of radical dependence
on explicit polynomial-ring modules with matching and with distinct
associated primes, verify total annihilator-dependence and its necessity
under a dimension-distinctness hypothesis, construct via localization a
module whose Radical Distinction Set is empty by design, confirm the
universal radical sum equivalence across several submodule pairs, and
check the multiplication-module identity exhaustively on the full
submodule lattice of $\mathbb{Z}/25\mathbb{Z}$. In every case the
computed annihilators and radicals agree exactly with the theoretical
predictions.
\end{abstract}

\textbf{Keywords:} annihilator ideals, radical submodules, associated
primes, multiplication modules, computer algebra, OSCAR, Julia.

\textbf{2020 Mathematics Subject Classification:} 13C13, 13C05, 13P99.

\section{Introduction}

In \cite{pekin-ozkaya-2026}, the authors introduce a notion of radical
dependence for submodules $N_1, N_2$ of a module $M$ over a commutative
Noetherian ring $R$: $N_1$ and $N_2$ are \emph{radically dependent} if
\[
  \sqrt{\operatorname{Ann}(N_1+N_2)} = \sqrt{\operatorname{Ann}(N_1)} +
  \sqrt{\operatorname{Ann}(N_2)},
\]
and it is shown (Theorem~1 of \cite{pekin-ozkaya-2026}) that this holds if
and only if
\[
  \sqrt{\operatorname{Ann}(N_1)} = \sqrt{\operatorname{Ann}(N_2)}.
\]

The purpose of this short note is purely computational: we use the
\textsc{OSCAR} system~\cite{oscar}, written in Julia and combining
\textsc{Singular}, \textsc{GAP}, \textsc{Polymake} and \textsc{Antic}, to
construct explicit modules and submodules, compute their annihilators and
radicals via Gr\"obner basis and primary decomposition algorithms, and
check the identity above symbolically rather than by hand. All
computations below were carried out with Julia~1.12.6 and
\texttt{Oscar.jl} version~1.8.0, as resolved from the General registry
in July~2026, and can be reproduced verbatim by pasting the listed code
into a Julia session after \texttt{using Oscar}. The exact dependency
graph (including the bundled \textsc{Singular}, \textsc{GAP},
\textsc{Polymake} and \textsc{Antic} components and their binary
artifacts) is pinned in the \texttt{Manifest.toml} file of the
accompanying repository, together with the complete, individually
runnable source code for every computation in this note.\footnote{\url{https://github.com/hamdooz/oscor-oriented-computation}}

\section{Elementary examples over $\mathbb{Z}$}

Examples~1--3 of \cite{pekin-ozkaya-2026} concern $M=\mathbb Z/n\mathbb Z$
as a $\mathbb Z$-module for $n=8,5,25$. Since $\mathbb Z$ is a principal
ideal domain, annihilators of cyclic submodules are computed directly
from $\gcd$, which we record here for completeness.

\begin{lstlisting}
n = 8
for m in 1:(n-1)
    a = n / gcd(m, n)
    println("Ann(", m, "-bar) = (", Int(a), ")Z")
end
\end{lstlisting}

This reproduces $\operatorname{Ann}(\bar 1)=\operatorname{Ann}(\bar
3)=\operatorname{Ann}(\bar 5)=\operatorname{Ann}(\bar 7)=8\mathbb Z$,
$\operatorname{Ann}(\bar 2)=\operatorname{Ann}(\bar 6)=4\mathbb Z$, and
$\operatorname{Ann}(\bar 4)=2\mathbb Z$, exactly as in Example~1 of
\cite{pekin-ozkaya-2026}. The same computation with $n=5$ and $n=25$
reproduces Examples~2 and~3 respectively.

\section{Verification of Theorem~1 over a polynomial ring}

For reference, we restate the result verbatim from \cite{pekin-ozkaya-2026}.

\begin{theorem}[Characterization of radical dependence; Theorem~1 of \cite{pekin-ozkaya-2026}]
\label{thm:1}
Let $N_1,N_2\le M$. Then $N_1$ and $N_2$ are radically dependent if and
only if $\sqrt{\Ann(N_1)}=\sqrt{\Ann(N_2)}$.
\end{theorem}

To test the characterization of radical dependence in a setting with a
non-principal-ideal-domain base ring, we work over $R=\mathbb Q[x,y]$ and
realize each submodule as an explicit \texttt{SubquoModule} of a rank-one
free module, following the construction used in the proof of Theorem~3 of
\cite{pekin-ozkaya-2026}: given a prime ideal $\mathfrak p\subset R$, the
cyclic module $R/\mathfrak p$ is represented as $F/\mathfrak p F$ for
$F\cong R$, and $M = R/\mathfrak p_1 \oplus R/\mathfrak p_2$ is formed as
a direct sum, with $N_1,N_2\le M$ the two canonical summands. For this
construction $\operatorname{Ann}(N_i)=\mathfrak p_i$ and
$\operatorname{Ann}(N_1+N_2) = \operatorname{Ann}(N_1)\cap
\operatorname{Ann}(N_2)$.

\subsection{Case $\mathfrak p_1 \neq \mathfrak p_2$}

Take $\mathfrak p_1 = (x)$ and $\mathfrak p_2 = (x-1,y)$, two distinct
prime ideals of $R=\mathbb Q[x,y]$.

\begin{lstlisting}
R, (x, y) = polynomial_ring(QQ, [:x, :y])

p1 = ideal(R, [x])
p2 = ideal(R, [x - 1, y])

F1 = free_module(R, 1)
N1_alone = SubquoModule(F1, R[one(R);], R[x;])

F2 = free_module(R, 1)
N2_alone = SubquoModule(F2, R[one(R);], R[x-1; y])

M, incs = direct_sum(N1_alone, N2_alone)
N1, _ = image(incs[1])
N2, _ = image(incs[2])

J1    = annihilator(N1)
J2    = annihilator(N2)
Jsum  = annihilator(N1 + N2)

rad_J1   = radical(J1)
rad_J2   = radical(J2)
rad_Jsum = radical(Jsum)

println(rad_J1 == p1)                              # true
println(rad_J2 == p2)                              # true
println(rad_Jsum == intersect(rad_J1, rad_J2))      # true
println(rad_Jsum == rad_J1 + rad_J2)                # false
\end{lstlisting}

The computation returns \texttt{rad\_J1 == p1} and \texttt{rad\_J2 ==
p2}, confirming that the constructed submodules indeed have the intended
annihilators. It further confirms
$\sqrt{\operatorname{Ann}(N_1+N_2)}=\sqrt{\operatorname{Ann}(N_1)}\cap
\sqrt{\operatorname{Ann}(N_2)}$ (always true, by the identity
$\operatorname{Ann}(N_1+N_2)=\operatorname{Ann}(N_1)\cap
\operatorname{Ann}(N_2)$ used in the proof of Theorem~1), while
$\sqrt{\operatorname{Ann}(N_1+N_2)} \neq
\sqrt{\operatorname{Ann}(N_1)}+\sqrt{\operatorname{Ann}(N_2)}$: since
$\mathfrak p_1\neq \mathfrak p_2$ are distinct maximal-type primes with
$\mathfrak p_1+\mathfrak p_2 = R$ (their associated varieties, the line
$x=0$ and the point $(1,0)$, are disjoint), the sum $\sqrt{J_1}+\sqrt{J_2}$
is the unit ideal, while $\sqrt{J_{\mathrm{sum}}}$ is a proper ideal. This
is precisely the failure of radical dependence predicted by Theorem~1
when $\sqrt{\operatorname{Ann}(N_1)}\neq \sqrt{\operatorname{Ann}(N_2)}$.

\subsection{Case $\mathfrak p_1 = \mathfrak p_2$}

Repeating the construction with $\mathfrak p_1=\mathfrak p_2=(x)$:

\begin{lstlisting}
p1 = ideal(R, [x])
p2 = ideal(R, [x])

F1 = free_module(R, 1)
N1_alone = SubquoModule(F1, R[one(R);], R[x;])
F2 = free_module(R, 1)
N2_alone = SubquoModule(F2, R[one(R);], R[x;])

M, incs = direct_sum(N1_alone, N2_alone)
N1, _ = image(incs[1])
N2, _ = image(incs[2])

J1   = annihilator(N1);   J2   = annihilator(N2)
Jsum = annihilator(N1 + N2)

rad_J1   = radical(J1)
rad_J2   = radical(J2)
rad_Jsum = radical(Jsum)

println(rad_Jsum == rad_J1 + rad_J2)   # true
\end{lstlisting}

Here the output is \texttt{true}: with $\mathfrak p_1=\mathfrak p_2$, the
radical dependence identity of Theorem~1 holds exactly, as predicted.

\section{Verification of Theorem~2 (total annihilator-dependence)}
\label{sec:thm2}

\begin{theorem}[Theorem~2 of \cite{pekin-ozkaya-2026}]
\label{thm:2}
Let $R$ be a commutative Noetherian ring with identity, and let $M$ be a
finitely generated $R$-module. If $\Ass(M)=\{\mathfrak p\}$, then $M$ is
totally annihilator-dependent, i.e.\ for all nonzero submodules
$N_1,N_2\le M$ with $\Ann(N_i)\neq 0$,
\[
  \dim R/\Ann(N_1+N_2) = \min\{\dim R/\Ann(N_1),\;
  \dim R/\Ann(N_2)\}.
\]
\end{theorem}

To test this on a module whose associated submodules genuinely have
\emph{different} annihilators (so that the identity is not trivially true
by equality of ideals), we take $R=\mathbb Q[x,y]$ and
$M=R/(x^2)$, whose defining ideal $(x^2)$ is already primary with
associated prime $(x)$; hence $\operatorname{Ass}(M)=\{(x)\}$ is a
singleton, as required by the hypothesis. Inside $M$ we consider the
cyclic submodules $N_1=R\bar x$ and $N_2=R\bar y$ generated by the images
of $x$ and $y$ respectively.

In OSCAR, the associated primes of an ideal are extracted from
\texttt{primary\_decomposition}, which returns a vector of
(primary, prime) pairs; there is no ideal-level
\texttt{associated\_primes} function in the version used here.

\begin{lstlisting}
R, (x, y) = polynomial_ring(QQ, [:x, :y])

I = ideal(R, [x^2])

pd = primary_decomposition(I)
println("Ass(R/I) = ", [P for (Q, P) in pd])   # {(x)}, singleton

F = free_module(R, 1)
Mmod = SubquoModule(F, R[one(R);], R[x^2;])   # M = R/(x^2)
g = gens(Mmod)[1]

N1, _ = sub(Mmod, [x*g])   # R * xbar
N2, _ = sub(Mmod, [y*g])   # R * ybar

J1   = annihilator(N1)
J2   = annihilator(N2)
Jsum = annihilator(N1 + N2)

println("J1 = ", J1, "   radical = ", radical(J1))
println("J2 = ", J2, "   radical = ", radical(J2))
println("Jsum = ", Jsum, " radical = ", radical(Jsum))

A1,   _ = quo(R, J1)
A2,   _ = quo(R, J2)
Asum, _ = quo(R, Jsum)

d1, d2, dsum = dim(A1), dim(A2), dim(Asum)
println("dim R/J1 = ", d1, ", dim R/J2 = ", d2, ", dim R/Jsum = ", dsum)
println("Theorem 2 holds? ", dsum == min(d1, d2))
\end{lstlisting}

The computation confirms \texttt{Ass(R/I) = \{(x)\}}, verifying the
hypothesis rather than merely assuming it. It further returns
$J_1=\operatorname{Ann}(N_1)=(x)$ and
$J_2=\operatorname{Ann}(N_2)=(x^2)$: these two submodules have
\emph{genuinely different} annihilators, so the equality
$\dim R/J_{\mathrm{sum}}=\min\{\dim R/J_1,\dim R/J_2\}$ is not a
tautology. Since $\sqrt{J_1}=\sqrt{J_2}=(x)$, both quotient rings have
Krull dimension $1$, and OSCAR confirms $\dim R/J_{\mathrm{sum}}=1=\min\{1,1\}$
as well, so the predicted equality holds exactly.

\section{Verification of Theorem~3 (necessity of a singleton $\Ass(M)$)}

\begin{theorem}[Theorem~3 of \cite{pekin-ozkaya-2026}]
\label{thm:3}
Let $R$ be a commutative Noetherian ring with identity, and let $M$ be a
finitely generated $R$-module. Suppose that for $\mathfrak p_1,\mathfrak
p_2\in\Ass(M)$, whenever $\mathfrak p_1\neq\mathfrak p_2$, one has
$\dim R/\mathfrak p_1 \neq \dim R/\mathfrak p_2$. If $M$ is totally
annihilator-dependent, then $\Ass(M)$ is a singleton.
\end{theorem}

We verify the contrapositive: we exhibit a module $M$ with
$\Ass(M)=\{\mathfrak p_1,\mathfrak p_2\}$, $\mathfrak p_1\neq \mathfrak
p_2$, satisfying the dimension-distinctness hypothesis, and show
computationally that $M$ is \emph{not} totally annihilator-dependent, by
exhibiting a witnessing pair $N_1,N_2$ for which the defining identity of
Definition~2 of \cite{pekin-ozkaya-2026} fails.

We reuse the construction of Section~3.1: $\mathfrak p_1=(x)$ (a line,
$\dim R/\mathfrak p_1=1$) and $\mathfrak p_2=(x-1,y)$ (a point, $\dim
R/\mathfrak p_2=0$), so the hypothesis $\dim R/\mathfrak p_1\neq\dim
R/\mathfrak p_2$ holds, and $M=R/\mathfrak p_1\oplus R/\mathfrak p_2$ has
$\Ass(M)=\{\mathfrak p_1,\mathfrak p_2\}$ via the two summands $N_1,N_2$.

\begin{lstlisting}
R, (x, y) = polynomial_ring(QQ, [:x, :y])

p1 = ideal(R, [x])
p2 = ideal(R, [x - 1, y])

F1 = free_module(R, 1)
N1_alone = SubquoModule(F1, R[one(R);], R[x;])
F2 = free_module(R, 1)
N2_alone = SubquoModule(F2, R[one(R);], R[x-1; y])

M, incs = direct_sum(N1_alone, N2_alone)
N1, _ = image(incs[1])
N2, _ = image(incs[2])

J1   = annihilator(N1)
J2   = annihilator(N2)
Jsum = annihilator(N1 + N2)

A1,   _ = quo(R, J1)
A2,   _ = quo(R, J2)
Asum, _ = quo(R, Jsum)

d1, d2, dsum = dim(A1), dim(A2), dim(Asum)

println("dim R/p1 = ", d1)
println("dim R/p2 = ", d2)
println("dim R/(p1 cap p2) = ", dsum)
println("Hypothesis of Theorem 3 (d1 != d2)? ", d1 != d2)
println("dsum == max(d1,d2)? ", dsum == max(d1, d2))
println("dsum == min(d1,d2)? ", dsum == min(d1, d2))
\end{lstlisting}

OSCAR returns $d_1=1$, $d_2=0$, and $d_{\mathrm{sum}}=1$. Thus
$d_{\mathrm{sum}}=\max\{d_1,d_2\}$, confirming the standard fact that
$\dim R/(\mathfrak p_1\cap\mathfrak p_2)=\max\{\dim R/\mathfrak
p_1,\dim R/\mathfrak p_2\}$ used in the proof of Theorem~3 in
\cite{pekin-ozkaya-2026}, and \emph{not} $\min\{d_1,d_2\}$. This is
exactly the failure of the total annihilator-dependence identity: since
$d_1\neq d_2$ (the theorem's hypothesis is met) and $\Ass(M)$ has two
elements, $M$ fails to be totally annihilator-dependent, precisely as the
contrapositive of Theorem~3 requires. Equivalently, this exhibits an
explicit obstruction: a singleton $\Ass(M)$ is not merely sufficient
(Theorem~2) but, under the dimension-distinctness hypothesis, necessary
for total annihilator-dependence.

\section{Verification of Theorem~4 (the Radical Distinction Set)}

\begin{definition}[Radical Distinction Set; Definition~3 of \cite{pekin-ozkaya-2026}]
Let $M$ be an $R$-module. The Radical Distinction Set of $M$ is
\[
  Z_g(M) = \{\, m\in M\setminus\{0\} \mid \Ann(Rm) \neq \Ann(\Rad(M))\,\}.
\]
\end{definition}

\begin{theorem}[Theorem~4 of \cite{pekin-ozkaya-2026}]
\label{thm:4}
Let $R$ be a Noetherian commutative ring and $M\neq 0$ a finitely
generated $R$-module. If $Z_g(M)=\emptyset$, then
$\Ass(M)=\{\Ann(\Rad(M))\}$, and $\Ann(\Rad(M))$ is a prime ideal.
\end{theorem}

OSCAR has no built-in \texttt{Rad} or \texttt{Zg} function, so both must
be assembled from primitives. A first obstacle is that neither of the
elementary examples of \cite{pekin-ozkaya-2026} (Examples~1 and~2,
i.e.\ $\mathbb Z/8\mathbb Z$ and $\mathbb Z/5\mathbb Z$) actually
satisfies the hypothesis $Z_g(M)=\emptyset$: in both cases $Z_g(M)\neq
\emptyset$, so Theorem~\ref{thm:4} is never triggered by the paper's own
illustrative examples. To exercise it meaningfully we construct a fresh
module for which $Z_g(M)=\emptyset$ can be established outright, rather
than merely assumed.

Take $R_0=\mathbb Q[x]$ and localize at the maximal ideal of the origin,
$U=R_0\setminus (x)$, to obtain the local domain $R=R_{(x)}$ with unique
maximal ideal $\mathfrak m=(x)R$. Let $M=R$, viewed as a module over
itself. Since $R$ is local, $\Rad(M)=\mathfrak m$ (the unique maximal
submodule). Since $R$ is furthermore a \emph{domain}, multiplication by
any nonzero ring element is injective, so $\Ann(Rm)=(0)$ for every
nonzero $m\in M$, and likewise $\Ann(\mathfrak m)=(0)$ (only $0$ can kill
a nonzero ideal of a domain). Hence \emph{every} nonzero $m$ satisfies
$\Ann(Rm)=\Ann(\Rad(M))=(0)$, so $Z_g(M)=\emptyset$ by a structural
argument -- not by an exhaustive search, which would be impossible here
since $R$ is infinite.

\begin{lstlisting}
R0, (x,) = polynomial_ring(QQ, [:x])
U = complement_of_point_ideal(R0, [0])
RL, iota = localization(R0, U)

xL = iota(x)

F = free_module(RL, 1)
Nm, _ = sub(F, [xL*F[1]])      # the maximal ideal m as a submodule of M = RL
Ann_m = annihilator(Nm)
println("Ann(Rad(M)) = ", Ann_m)

zero_ideal = ideal(RL, [zero(RL)])
pd = primary_decomposition(zero_ideal)
println("Ass(M) = ", [P for (Q, P) in pd])

println("Ann(Rad(M)) == (0)? ", Ann_m == zero_ideal)
println("Theorem 4 holds? Ass(M) = {Ann(Rad(M))}? ",
        length(pd) == 1 && pd[1][2] == Ann_m)
\end{lstlisting}

OSCAR confirms $\Ann(\Rad(M))=(0)$, and \texttt{primary\_decomposition}
of the zero ideal returns the single pair $((0),(0))$, so
$\Ass(M)=\{(0)\}$. Both final checks evaluate to \texttt{true}: the
predicted equality $\Ass(M)=\{\Ann(\Rad(M))\}$ holds exactly, and the
ideal $(0)$ is indeed prime, as required by Theorem~\ref{thm:4}.

\begin{remark}
This example is complementary to the module $M=R/(x^2)$ of
Section~4: that module also has a singleton $\Ass(M)=\{(x)\}$, yet
$Z_g(M)\neq\emptyset$ there (the submodules $R\bar x$ and $R\bar y$ have
different annihilators $(x)\neq(x^2)$, neither equal to $\Ann(\Rad(M))$
in general). This shows that the hypothesis of Theorem~\ref{thm:4} is
sufficient, but not necessary, for $\Ass(M)$ to be a singleton --
consistent with Theorem~\ref{thm:2} providing a strictly weaker
sufficient condition (singleton $\Ass(M)$ alone) for the related but
distinct conclusion of total annihilator-dependence.
\end{remark}

\section{Verification of Theorem~5 (radical sum equivalence)}
\label{sec:thm5}

\begin{theorem}[Theorem~5 of \cite{pekin-ozkaya-2026}]
\label{thm:5}
Let $R$ be a Noetherian commutative ring and $M$ a finitely generated
$R$-module. The following are equivalent:
\begin{enumerate}
\item[(i)] for all nonzero $N_1,N_2\le M$: $\sqrt{\Ann(N_1+N_2)} =
\sqrt{\Ann(N_1)}+\sqrt{\Ann(N_2)}$;
\item[(ii)] $\Ass(M)=\{\mathfrak p\}$ for some prime ideal $\mathfrak p$.
\end{enumerate}
\end{theorem}

Theorem~\ref{thm:5} strengthens Theorem~\ref{thm:1} from a statement about
a \emph{single} pair $N_1,N_2$ to a statement about \emph{all} pairs
simultaneously. The direction (i)$\Rightarrow$(ii) is exactly the
contrapositive already exhibited in Section~3.1: there, with $\Ass(M)$
containing the two distinct primes $\mathfrak p_1=(x)$ and
$\mathfrak p_2=(x-1,y)$, a single witnessing pair $N_1,N_2$ was already
enough to falsify the radical sum identity, confirming $\neg$(i) whenever
$\neg$(ii) holds -- no further computation is needed for that direction.

For (ii)$\Rightarrow$(i), we reuse the module $M=R/(x^2)$ of Section~4,
whose associated primes were confirmed to be the singleton $\{(x)\}$ via
\texttt{primary\_decomposition}. Since a single instance cannot establish
a claim about \emph{all} pairs $N_1,N_2$, we test the identity on four
structurally distinct pairs of cyclic submodules, including pairs whose
annihilators $\Ann(N_1)$ and $\Ann(N_2)$ are themselves unequal.

\begin{lstlisting}
R, (x, y) = polynomial_ring(QQ, [:x, :y])
I = ideal(R, [x^2])

pd = primary_decomposition(I)
println("Ass(R/I) = ", [P for (Q, P) in pd])   # {(x)}

F = free_module(R, 1)
Mmod = SubquoModule(F, R[one(R);], R[x^2;])
g = gens(Mmod)[1]

pairs = [
    (x*g, y*g),
    (y*g, (x+y)*g),
    ((2*x+3*y)*g, (x-y)*g),
    (x*g, (x+y)*g),
]

for (i, (a, b)) in enumerate(pairs)
    Na, _ = sub(Mmod, [a])
    Nb, _ = sub(Mmod, [b])
    Ja   = annihilator(Na)
    Jb   = annihilator(Nb)
    Jsum = annihilator(Na + Nb)
    lhs = radical(Jsum)
    rhs = radical(Ja) + radical(Jb)
    println("Pair $i: J_a=", Ja, "  J_b=", Jb, "  lhs==rhs? ", lhs == rhs)
end
\end{lstlisting}

All four pairs return \texttt{lhs==rhs? true}, including Pairs~1 and~4,
in which $\Ann(N_1)=(x)$ and $\Ann(N_2)=(x^2)$ are themselves
\emph{unequal} ideals: the radical sum identity still holds because both
annihilators share the same radical $(x)=\sqrt{\Ann(M)}$, consistent with
Theorem~\ref{thm:1}. While a finite sample of pairs cannot constitute a
proof of the universal statement (ii)$\Rightarrow$(i) for all pairs over
an infinite ring, the consistency across four structurally different
choices -- together with the structural argument already available from
Theorem~\ref{thm:2} (that $\Ass(M)=\{\mathfrak p\}$ forces
$\sqrt{\Ann(N)}=\mathfrak p$ for every nonzero submodule $N\le M$, from
which (i) follows immediately) -- corroborates the equivalence.

\section{Verification of Theorem~6 (multiplication modules)}

\begin{theorem}[Theorem~6 of \cite{pekin-ozkaya-2026}]
\label{thm:6}
Let $R$ be a Noetherian commutative ring and $M$ a finitely generated
multiplication $R$-module with $\Ass(M)=\{\mathfrak p\}$. Then for all
nonzero $N_1,N_2\le M$,
\[
  \sqrt{\Ann(N_1+N_2)} = \sqrt{\Ann(N_1)} + \sqrt{\Ann(N_2)} = \mathfrak p.
\]
\end{theorem}

Recall that an $R$-module $M$ is a multiplication module if every
submodule of $M$ has the form $IM$ for some ideal $I\le R$. Every cyclic
module $M=R/I$ is trivially a multiplication module, since its
submodules are exactly $J/I$ for ideals $J\ge I$, i.e.\ $(J/I)\cdot(R/I)$.
In particular, the module $M=R/(x^2)$ already used in
Sections~\ref{sec:thm2} and~\ref{sec:thm5} is a multiplication module, and the
computations recorded there -- in which several pairs of cyclic
submodules were shown to satisfy
$\sqrt{\Ann(N_1+N_2)}=\sqrt{\Ann(N_1)}+\sqrt{\Ann(N_2)}=(x)$ -- already
constitute a verification of Theorem~\ref{thm:6} in the polynomial-ring
setting, with the common value equal to the unique associated prime
$(x)$, exactly as the theorem asserts.

To complement this with the ring-theoretic setting of Example~3 of
\cite{pekin-ozkaya-2026}, we verify the identity directly and
exhaustively on $M=\mathbb Z/25\mathbb Z$, whose complete submodule
lattice is the chain $0 \subset 5\mathbb Z/25\mathbb Z \subset M$, hence
small enough to check every pair of nonzero submodules explicitly.

\begin{lstlisting}
n = 25
p = 5
submods = [5, 1]   # generators of 5Z/25Z and of M = Z/25Z itself

for a in submods, b in submods
    ann_a    = div(n, gcd(a, n))
    ann_b    = div(n, gcd(b, n))
    gen_sum  = gcd(a, b)
    ann_sum  = div(n, gcd(gen_sum, n))

    rad_is_p = (ann_sum % p == 0) && (ann_a % p == 0) && (ann_b % p == 0)
    println("Ann(N1)=", ann_a, "Z, Ann(N2)=", ann_b,
            "Z, Ann(N1+N2)=", ann_sum,
            "Z -- radical of all three equals p=", p, "Z? ", rad_is_p)
end
\end{lstlisting}

All four pairs $(a,b)\in\{5,1\}^2$ (representing the two nonzero
submodules $5\mathbb Z/25\mathbb Z$ and $M$ itself, and their mutual
combinations) return \texttt{true}: the annihilators involved are always
$5\mathbb Z$ or $25\mathbb Z$, both of which have radical $5\mathbb Z=
\mathfrak p$, so the common-value identity of Theorem~\ref{thm:6} holds
throughout the entire submodule lattice of $M$, confirming
$\operatorname{Ass}(M)=\{5\mathbb Z\}$ from Example~3 of
\cite{pekin-ozkaya-2026} directly at the level of every possible pair of
submodules, rather than merely for the single generating example given
there.

\section{Conclusion}

The computations above instantiate both directions of the equivalence in
Theorem~1 of \cite{pekin-ozkaya-2026}, the conclusion of Theorem~2, the
contrapositive of Theorem~3, the hypothesis-and-conclusion pair of
Theorem~4, the universal radical sum equivalence of Theorem~5, and the
multiplication-module identity of Theorem~6, on genuinely non-principal
polynomial, localized-ring, and cyclic integer-quotient data,
independently of the hand computations given in the original proofs.
Together with the direct reproduction of the elementary $\mathbb
Z/n\mathbb Z$ examples, this provides a symbolic, machine-checked
consistency check of all six of the paper's central characterizations.
All source
code is included above in full so that the computations can be
reproduced by any reader with a working OSCAR installation.

\begin{thebibliography}{9}

\bibitem{pekin-ozkaya-2026}
A.~Pekin and H.~\"Ozkaya,
\emph{Annihilator-based dependency relations in modules and radical
characterizations},
Algebra and Discrete Mathematics \textbf{41} (2026), no.~2, 235--243.
\url{https://doi.org/10.12958/adm2471}

\bibitem{oscar}
The OSCAR Team,
\emph{OSCAR -- Open Source Computer Algebra Research system},
Version~1.x, 2026.
\url{https://www.oscar-system.org}

\end{thebibliography}

\end{document}
